\section{선별 문제 풀이 I}
\label{sec:equation-groups-solutions-a}

\subsection*{문제 \thechapter.1: 이차다항식 세 개}

\begin{solution}
(a) $X^2-5$의 근은 $\pm\sqrt5$이고 분해체는 $\Q(\sqrt5)$이다. $5$는 유리수의 제곱이 아니므로 차수는 $2$이고
\[
  \Gal(f/\Q)=\{\id,\sigma\}\cong C_2,
  \qquad
  \sigma(\sqrt5)=-\sqrt5.
\]
근에 대한 작용은 두 근을 바꾸는 전치이다.

(b) $X^2+4X+1$의 판별량은
\[
  D=4^2-4=12.
\]
근은 $-2\pm\sqrt3$이고 분해체는 $\Q(\sqrt3)$이다. 따라서 갈루아군은 $C_2$이며 비항등원은 두 근을 바꾼다.

(c) $X^2-4=(X-2)(X+2)$는 이미 $\Q$에서 완전히 분해된다. 분해체가 $\Q$이므로 갈루아군은 자명군이다.
\end{solution}

\subsection*{문제 \thechapter.2: 세 기약 삼차식}

\begin{solution}
(a) $X^3-2$는 $2$에 대한 아이젠슈타인 판정법으로 기약이다. 우울형에서 $a=0$, $b=-2$이므로
\[
  \Delta=-4a^3-27b^2=-108.
\]
이는 $\Q$의 제곱이 아니므로 갈루아군은 $S_3$이다.

(b) $X^3-X+1$은 가능한 유리근 $\pm1$을 갖지 않아 기약이다. 판별식은
\[
  \Delta=-4(-1)^3-27=4-27=-23
\]
이고 제곱이 아니므로 갈루아군은 $S_3$이다.

(c) $X^3-3X+1$도 가능한 유리근 $\pm1$을 갖지 않아 기약이다. 판별식은
\[
  \Delta=-4(-3)^3-27=108-27=81=9^2.
\]
따라서 갈루아군은 $A_3\cong C_3$이다.
\end{solution}

\subsection*{문제 \thechapter.3: reducible cubic의 궤도}

\begin{solution}
(a)
\[
  X^3-4X=X(X-2)(X+2)
\]
는 $\Q$에서 완전히 분해된다. 분해체는 $\Q$, 갈루아군은 자명군이고 세 궤도는 각각 한 원소로 이루어진다.

(b) $(X-1)(X^2-2)$의 분해체는 $\Q(\sqrt2)$이다. 갈루아군은 $C_2$이고 근 집합의 궤도는
\[
  \{1\},
  \qquad
  \{\sqrt2,-\sqrt2\}
\]
이다.

(c) 세 근이 모두 유리수이므로 분해체는 $\Q$, 갈루아군은 자명군이고 세 개의 한 원소 궤도가 있다.
\end{solution}

\subsection*{문제 \thechapter.5: 이차항이 있는 삼차식}

\begin{solution}
$f(X)=X^3+3X^2-1$에서 $X=Y-1$을 대입하면
\[
\begin{aligned}
  f(Y-1)
  &=(Y-1)^3+3(Y-1)^2-1\\
  &=Y^3-3Y+1.
\end{aligned}
\]
원래 다항식은 가능한 유리근 $\pm1$을 갖지 않으므로 기약이다. 평행이동은 분해체와 갈루아군을 바꾸지 않는다. 우울형의 판별식은
\[
  \Delta=-4(-3)^3-27=81=9^2.
\]
따라서
\[
  \Gal(f/\Q)\cong C_3.
\]
\end{solution}

\subsection*{문제 \thechapter.7: 사차 순수방정식}

\begin{solution}
$\alpha=\sqrt[4]{2}$라 하자. $X^4-2$의 근은
\[
  \alpha,-\alpha,i\alpha,-i\alpha
\]
이므로 분해체는 $L=\Q(\alpha,i)$이다. 아이젠슈타인 판정법으로
\[
  [\Q(\alpha):\Q]=4.
\]
$\Q(\alpha)$는 실수체이고 $i$는 비실수이므로 $i\notin\Q(\alpha)$이다. 따라서
\[
  [L:\Q]=2\cdot4=8.
\]

주어진 사상은 생성원들이 만족하는 관계 $\alpha^4=2$, $i^2=-1$을 보존하고 분해체 안에서 근들을 순열하므로 자동동형을 정의한다. 직접 계산하면
\[
  r^4=s^2=1
\]
이고
\[
  srs(\alpha)=s r(\alpha)=s(i\alpha)=-i\alpha=r^{-1}(\alpha),
  \qquad
  srs(i)=i=r^{-1}(i).
\]
따라서 $srs=r^{-1}$이다. 여덟 원소
\[
  1,r,r^2,r^3,s,sr,sr^2,sr^3
\]
가 서로 다르고 갈루아군 위수도 $8$이므로
\[
  \Gal(L/\Q)\cong D_4.
\]
\end{solution}

\subsection*{문제 \thechapter.10: 이차식과 삼차식의 합성 분해체}

\begin{solution}
$E=\Q(\sqrt2)$를 $X^2-2$의 분해체, $M$을 $X^3-3$의 분해체라 하자. $X^3-3$은 기약이고 판별식은
\[
  -27\cdot 3^2=-243
\]
이므로
\[
  \Gal(M/\Q)\cong S_3,
  \qquad [M:\Q]=6.
\]
$M$의 유일한 이차 중간체는
\[
  \Q(\sqrt{-243})=\Q(\sqrt{-3}).
\]
따라서 $E\ne\Q(\sqrt{-3})$이고
\[
  E\cap M=\Q.
\]
전체 다항식의 분해체는 $L=EM$이며
\[
  [L:\Q]
  =\frac{[E:\Q][M:\Q]}{[E\cap M:\Q]}
  =2\cdot6=12.
\]
두 갈루아 확장이 교집합 $\Q$을 가지므로 제한 사상으로
\[
  \Gal(L/\Q)
  \cong\Gal(E/\Q)\times\Gal(M/\Q)
  \cong C_2\times S_3.
\]
근 집합의 궤도는 이차인수의 두 근과 삼차인수의 세 근으로 나뉜다.
\end{solution}

\subsection*{문제 \thechapter.12: 기약성과 추이성}

\begin{solution}
$f$가 기약이고 $\alpha,\beta$가 두 근이라 하자. 두 근의 최소다항식이 같으므로 $\alpha\mapsto\beta$인 $K$-동형
\[
  K(\alpha)\longrightarrow K(\beta)
\]
가 존재한다. 이를 대수적 폐포의 $K$-자동동형으로 연장하고 분해체 $L$에 제한한다. $L/K$가 정규이므로 제한은 $L$의 $K$-자동동형이며 $\alpha$를 $\beta$로 보낸다. 따라서 작용은 추이적이다.

반대로 $f=gh$가 두 비상수 $K$-다항식의 곱이라고 하자. 분리성 때문에 $g$와 $h$의 근 집합은 겹치지 않는다. 모든 $K$-자동동형은 $g$의 근 집합을 보존하므로 전체 근 집합에 비자명한 불변부분집합이 생긴다. 따라서 작용은 추이적이지 않다.
\end{solution}

\subsection*{문제 \thechapter.14: 소수차에서의 긴 순환}

\begin{solution}
$f$가 기약이므로 갈루아군 $G$는 $p$개의 근 위에서 추이적으로 작용한다. 한 근 $\alpha$에 대해 궤도-안정자 정리로
\[
  p=|G\cdot\alpha|=[G:\Stab_G(\alpha)].
\]
따라서 $p\mid|G|$이다. 코시 정리에 의해 $G$에는 위수 $p$인 원소 $\sigma$가 있다. $G\le S_p$에서 위수 $p$인 순열의 순환분해는 길이 $p$인 순환 하나와 고정점들로 이루어져야 한다. 움직일 수 있는 점이 정확히 $p$개이므로 $\sigma$는 $p$-순환이다.
\end{solution}
